%File: aaai2026-unified-supp.tex
%
% UNIFIED AAAI 2026 SUPPLEMENTARY MATERIAL TEMPLATE
% To switch between anonymous submission and camera-ready versions,
% simply change the next line:
%
% For ANONYMOUS SUBMISSION: uncomment the next line
\def\aaaianonymous{true}
%
% For CAMERA-READY VERSION: comment out or delete the next line
% \def\aaaianonymous{true}
%
%%%%%%%%%%%%%%%%%%%%%%%%%%%%%%%%%%%%%%%%%%%%%%%%%%%%%%%%%%%%%%%%%%%%%%%
%
% NOTE (2026-07-20): this file documents the method and BOTH experimental venues of the main
% paper. Section A specifies the METHOD. Section B is the SIMULATION venue
% (residual RL trained directly against the DexMimicGen/LIBERO simulators); Section C
% is the REAL-ROBOT venue (the shared residual objective trained inside a learned
% dynamics model). It remains a protocol-forward draft: every quantity here is
% either verified from the actual pipeline or marked \todo{}. The main paper carries
% the compact headline tables; this supplement reserves reader-facing per-seed and
% paired-result sections, which remain TODO while the simulation inventory is reconciled
% and the real-robot evaluation is pending. The simulation hyperparameters in
% Section B.3 were extracted from the trainer/config source
% (resfit/rl_finetuning/); the naming convention that encodes them in the run
% scripts is: as005 = action_scale 0.05, bc01 = demo_bc_coef 0.1,
% hiqlv512 = goal-conditioned value width 512, sg15 = subgoal_way_steps 15.
%
%%%%%%%%%%%%%%%%%%%%%%%%%%%%%%%%%%%%%%%%%%%%%%%%%%%%%%%%%%%%%%%%%%%%%%%

\documentclass[letterpaper]{article} % DO NOT CHANGE THIS

% Conditional package loading based on version
\ifdefined\aaaianonymous
    \usepackage[submission]{aaai2026}  % Anonymous submission version
\else
    \usepackage{aaai2026}              % Camera-ready version
\fi

\usepackage{times}  % DO NOT CHANGE THIS
\usepackage{helvet}  % DO NOT CHANGE THIS
\usepackage{courier}  % DO NOT CHANGE THIS
\usepackage[hyphens]{url}  % DO NOT CHANGE THIS
\usepackage{graphicx} % DO NOT CHANGE THIS
\urlstyle{rm} % DO NOT CHANGE THIS
\def\UrlFont{\rm}  % DO NOT CHANGE THIS
\usepackage{natbib}  % DO NOT CHANGE THIS AND DO NOT ADD ANY OPTIONS TO IT
\usepackage{caption} % DO NOT CHANGE THIS AND DO NOT ADD ANY OPTIONS TO IT
\frenchspacing  % DO NOT CHANGE THIS
\setlength{\pdfpagewidth}{8.5in} % DO NOT CHANGE THIS
\setlength{\pdfpageheight}{11in} % DO NOT CHANGE THIS

\usepackage{amsmath}
\usepackage{amssymb}
\usepackage{booktabs}

\pdfinfo{
/TemplateVersion (2026.1)
}

\setcounter{secnumdepth}{0} %May be changed to 1 or 2 if section numbers are desired.

% draft-only helper (remove before submission), matching main.tex
\newcommand{\todo}[1]{\textbf{[TODO: #1]}}
\newcounter{algorithm}

% Title - conditionally set based on version
\ifdefined\aaaianonymous
    \title{SHORE-RL: Shortening the Horizon for Long-Horizon Residual\\
    Reinforcement Learning\\
    \vspace{0.5em}\large Supplementary Material}
\else
    \title{SHORE-RL: Shortening the Horizon for Long-Horizon Residual\\
    Reinforcement Learning\\
    \vspace{0.5em}\large Supplementary Material}
\fi

% Author and affiliation information
\ifdefined\aaaianonymous
\author{
    Anonymous Submission
}
\affiliations{
    % Leave affiliations empty for anonymous submission
}
\else
\author{
    Anonymous Submission
}
\affiliations{
}
\fi

\begin{document}

\maketitle

\begin{abstract}
This supplement provides the complete SHORE-RL objective and implementation details,
followed by the simulation and real-robot experimental protocols. The two venues share
the same residual policy parameterization and actor--critic update, while instantiating
different transition collectors, reward sources, and replay samplers. We additionally
document the ablation design, run provenance, dynamics-model validation, and paired
physical-robot evaluation needed to interpret the main-paper results.
\end{abstract}

% Supplementary Content Starts Here

\section{A\quad Method and Implementation Details}

The main paper introduces the residual controller, self-derived waypoint, and two
potential choices. Here we give the complete optimization objective and the venue-parameterized update of
Algorithm~\ref{alg:shore}, in the notation the two venues reuse. Venue-independent settings
are collected in Table~\ref{tab:sharedhp}; venue-specific numbers appear in Section~B.3 and
Section~C.5.

\subsection{A.1\quad Residual Coordinates and Decision Time Scale}

All policy/value inputs use the frozen base feature $s_t=\psi(o_t)$. The executed action is
\[
a_t=a_b+s_a a_{\mathrm{res}},\qquad
a_{\mathrm{res}}=\tanh\!\big(\mu_\theta(s_t,z_t)\big)\in[-1,1]^d,
\]
where $a_b\sim\pi_b(\cdot\mid o_t)$. The actor's final layer is zero-initialized, so the
initial policy exactly equals the frozen base; base action chunks are executed open-loop.
\todo{define the RL decision time scale precisely: state whether the residual is produced
once per base-policy chunk or once per low-level control step, how it is broadcast over a
chunk, and whether episode horizons count control steps or residual decisions.}

\subsection{A.2\quad Waypoint Implementation Details}

As described in the main paper, an expectile-trained goal-conditioned value $V_g(s,g)$ and
an advantage-weighted navigator $\pi_h$ are initialized from demonstrations and frozen
features. Future goals follow the geometric schedule in Table~\ref{tab:simvalue}; the
navigator emits the normalized bottleneck waypoint
$z_t=\phi([s_t,s_{t+k}])$, which conditions the residual actor and critic.
\todo{reconcile the implementation's $0/-1$ goal-reward convention with the main paper's
$1/0$ description and retain the convention used by the actual training code.}
During residual RL, $\phi$ remains frozen, while $V_g$ and $\pi_h$ are jointly updated in
the waypoint-enabled arms.

\subsection{A.3\quad Potential Implementations}

Both potential choices use
\[
\tilde r_t=r_t+\gamma\Phi(s_{t+1})-\Phi(s_t),
\]
with $\Phi$ frozen. The simulation headline arm uses a task-specific stage potential
$\Phi_{\mathrm{stage}}$.
\todo{confirm that the reported stage bonus is implemented as the static difference
$\gamma\Phi_{\mathrm{stage}}(s')-\Phi_{\mathrm{stage}}(s)$ rather than as a direct
transition bonus; otherwise describe the implemented reward and remove the invariance
claim.}
The privilege-free simulation arm and all imagination training instead use a separate
single-state value,
\[
\Phi_V(s)=cV(s),\qquad
c=\frac{N_{\mathrm{stage}}-1}{V_{\max}-V_{\min}},
\]
trained from demonstrations with reward $1$ at the final step and $0$ elsewhere.
\todo{define $c$ for tasks without a stage decomposition, including LIBERO and the
real-robot tasks.} The goal-conditioned $V_g$ supplies only the waypoint; the frozen
single-state $V$ supplies only $\Phi_V$.
\todo{confirm terminal-potential and truncation handling in both venues before claiming
potential-based policy invariance.}

\subsection{A.4\quad Residual Off-Policy Reinforcement Learning}

The residual is trained off-policy in the RLPD regime with TD3-style updates on a
mixed replay buffer. An ensemble of $N$ critics $\{Q_i\}_{i=1}^{N}$ (with $N=10$) is
trained toward an $n$-step clipped target that takes the minimum over two randomly chosen
heads,
\[
y_t=\sum_{j=0}^{n-1}\gamma^j\tilde r_{t+j}
+\gamma^n(1-d_{t:t+n})
\min_{i\in\mathcal{B}}\bar Q_i\big(s_{t+n},z_{t+n},a'_{\mathrm{res}}\big),
\]
where $d_{t:t+n}$ masks a terminal transition within the target window,
$\mathcal{B}$ is a pair of random heads, and
$a'_{\mathrm{res}}=\operatorname{clip}(\tanh(\mu_{\bar\theta})+\epsilon,-1,1)$ carries
clipped target-smoothing noise. The critic loss is
$\sum_i\mathbb{E}[(Q_i(s,z,a_{\mathrm{res}})-y_t)^2]$. The actor is updated less often
than the critic and anchored to the demonstrations,
\[
\mathcal{L}_{\mu}=-\,\mathbb{E}\big[\bar Q(s,z,a_{\mathrm{res}}^\theta)\big]
+\lambda_{\mathrm{bc}}\,\mathbb{E}_{\mathcal{D}}
\big[\lVert a_{\mathrm{res}}^\theta-a_{\mathrm{res}}^{\mathcal D}\rVert^2\big],
\]
where $a_{\mathrm{res}}^\theta=\tanh(\mu_\theta(s,z))$ and
$a_{\mathrm{res}}^{\mathcal D}$ denotes the demonstrator correction expressed in the
same coordinates as the actor output. For the normalized parameterization of
Section~A.1 this would ordinarily be
$\operatorname{clip}((a_{\mathrm{demo}}-a_b)/s_a,-1,1)$.
\todo{verify whether the implementation divides the demonstration correction by $s_a$,
and make the equation, replay contents, and Table~\ref{tab:simresrl} match the code.}
The residual actor is zero-initialized as in Section~A.1, and $s_a$ keeps every explored
action near the base.

\subsection{A.5\quad Shared Update, Two Venue Instantiations}

The actor and critic updates do not depend on where a transition originates, but the
collector $\mathcal C$, reward source $R$, potential $\Phi$, and replay sampler $\rho$ do.
Algorithm~\ref{alg:shore} therefore states a venue-parameterized update rather than
claiming that only one collection line changes. Table~\ref{tab:sharedhp} lists the
venue-independent hyperparameters both venues reuse, and Table~\ref{tab:venueinst} lists
what each venue sets differently.

\begin{figure}[t]
\small
\refstepcounter{algorithm}\label{alg:shore}
\setlength{\tabcolsep}{2pt}
\begin{tabular}{@{}r@{\,\,}p{0.84\linewidth}@{}}
\toprule
\multicolumn{2}{@{}p{0.9\linewidth}@{}}{\textbf{Algorithm \thealgorithm}\quad SHORE-RL residual training across validation venues}\\[1pt]
\multicolumn{2}{@{}p{0.9\linewidth}@{}}{\textbf{Input} frozen $\pi_b,\psi$; goal value $V_g$ and navigator $\pi_h$; frozen bottleneck $\phi$; potential $\Phi$; collector $\mathcal C$; reward source $R$; sampler $\rho$; demonstrations $\mathcal D$}\\[1pt]
\multicolumn{2}{@{}p{0.9\linewidth}@{}}{\textbf{Output} trained residual actor $\mu_\theta$}\\
\midrule
$1$ & initialize a zero-output residual actor and a $10$-head critic ensemble \\
$2$ & initialize replay stores with the offline transitions required by $\rho$ \\
$3$ & \textbf{for} each collection round \textbf{do} \\
$4$ & \quad $\mathcal T\gets\mathcal C(\pi_b,\mu_\theta,\pi_h,\psi)$ \\
$5$ & \quad \textbf{for} each $(s,z,a_{\mathrm{res}},r_{\mathrm{env}},s',d)\in\mathcal T$ \textbf{do} \\
$6$ & \quad\quad $\tilde r\gets R(r_{\mathrm{env}})+\gamma\Phi(s')-\Phi(s)$; store the transition \\
$7$ & \quad \textbf{for} each critic update prescribed by UTD \textbf{do} \\
$8$ & \quad\quad sample a venue-specific mixed batch using $\rho$ \\
$9$ & \quad\quad update all critics toward the $n$-step clipped target \\
$10$ & \quad\quad on delayed steps, update the actor with $Q$ and the demo anchor \\
$11$ & \quad\quad soft-update actor and critic targets \\
$12$ & \quad\quad update $V_g$ and $\pi_h$ on mixed data; keep $\phi$ frozen \\
$13$ & \textbf{return} $\mu_\theta$ \\
\bottomrule
\end{tabular}
\end{figure}

\begin{table}[t]
\centering\small\setlength{\tabcolsep}{3pt}
\begin{tabular}{@{}p{0.23\columnwidth}p{0.31\columnwidth}p{0.36\columnwidth}@{}}
\toprule
Component & Simulation & Real-robot imagination \\
\midrule
Collector $\mathcal C$ & simulator steps & bounded rollouts of frozen $D$ \\
Env. reward & task reward & none ($r_{\mathrm{env}}=0$) \\
Potential $\Phi$ & stage or $\Phi_V$ & $\Phi_V$ only \\
Replay mix $\rho$ & demo $+$ online & recorded $+$ imagined \\
Re-anchor & simulator reset & real-data start; depth $\leq2$ \\
Other params & Table~\ref{tab:simresrl} & Section~C.5 \\
\bottomrule
\end{tabular}
\caption{Shared update and venue-specific instantiations. The policy and critic objectives
are common; collection, reward availability, and replay composition are not. The exact
real/imagined/demo mixture remains to be reported in Section~C.5.}
\label{tab:venueinst}
\end{table}

\begin{table}[t]
\centering\small\setlength{\tabcolsep}{4pt}
\begin{tabular}{@{}ll@{}}
\toprule
Setting & Value \\
\midrule
Residual output / scale $s_a$   & $\tanh\times s_a$, $s_a=0.05$, zero-init final layer \\
Critic ensemble                 & $10$ Q-nets; target $\min$ of $2$ random heads \\
Critic / actor MLP              & $2$ layers $\times$ $1024$, ReLU, LayerNorm \\
Discount $\gamma$ / $n$-step    & $0.99$ / $3$ \\
UTD / target Polyak             & $4$ / $0.005$ \\
Target smoothing                & $0.05$, clip $0.3$ \\
Demo-BC coefficient $\lambda_{\mathrm{bc}}$ & $0.1$ \\
Waypoint bottleneck dim         & $10$ (radius $\sqrt{10}$) \\
Expectile $\tau$ / AWR $\beta$  & $0.7$ / $1.0$ \\
actor\_lr / critic\_lr          & $1{\times}10^{-6}$ / $1{\times}10^{-4}$ (AdamW) \\
\bottomrule
\end{tabular}
\caption{Venue-independent SHORE-RL hyperparameters, shared by the simulation venue
(Section~B) and the real-robot venue (Section~C). Venue-specific settings (budgets,
collection, replay composition, evaluation) appear in Table~\ref{tab:simresrl} and
Section~C.5.}
\label{tab:sharedhp}
\end{table}

\section{B\quad Simulation Venue}

In simulation the environment is cheap to step, so the residual is trained
\emph{directly} against the simulator with no learned dynamics model. This section
documents that venue end to end. The algorithm is the one specified in Section~A, a frozen base policy with a small
additive off-policy residual, a waypoint from a goal-conditioned value, and a dense
potential; this section fixes its concrete configuration for the simulator, and the
shared and venue-specific components are separated in Table~\ref{tab:venueinst}.

\subsection{B.1\quad Simulated Tasks and Success Criteria}

We evaluate on five DexMimicGen~\cite{jia2024dexmimicgen,mandlekar2023mimicgen}
dual-arm tasks, comprising four long-horizon (\textsc{ThreePieceAssembly}, \textsc{Pouring},
\textsc{LiftTray}, \textsc{Threading}) and the shorter-horizon \textsc{CanSort} as a
same-benchmark control, together with a single-arm extension on
LIBERO-10~\cite{liu2023libero}. Every DexMimicGen task is simulated at a $20$\,Hz
control rate; each episode is capped at the per-task horizon in
Table~\ref{tab:simtasks}. Success is the task's own simulator predicate evaluated at
the end of the episode, and that same predicate defines both the sparse terminal
reward and, on the stage-structured tasks, the per-stage bonus of the main result
(the stage count is the number of subtasks the detector recognizes; \textsc{CanSort}
has no stage decomposition and is run with no shaping).

\begin{table}[t]
\centering\small\setlength{\tabcolsep}{4pt}
\begin{tabular}{@{}lllrr@{}}
\toprule
Task & Robots & act.\ & horizon & stages \\
\midrule
\textsc{ThreePieceAssembly} & $2{\times}$Panda    & $14$ & $500$ & $4$ \\
\textsc{Pouring}            & GR1                 & $24$ & $400$ & $5$ \\
\textsc{LiftTray}           & $2{\times}$PandaDex  & $24$ & $650$ & $4$ \\
\textsc{Threading}          & $2{\times}$Panda    & $14$ & $300$ & $3$ \\
\textsc{CanSort} (control)  & GR1                 & $24$ & $400$ & none \\
\bottomrule
\end{tabular}
\caption{Simulated DexMimicGen tasks (all $20$\,Hz). \emph{act.}\ is the action
dimension; \emph{horizon} is the max episode length used by our rollout wrapper;
\emph{stages} is the number of subtasks in the stage decomposition (\textsc{CanSort}
has none). Success criteria (task predicate at episode end): \textsc{ThreePieceAssembly} requires both
pieces assembled onto the base and released; \textsc{Pouring} requires the ball transferred into
the bowl and the bowl placed upright on the pad; \textsc{LiftTray} requires both cubes on the
tray and the tray lifted clear of the table; \textsc{Threading} requires the needle inserted
through the ring; \textsc{CanSort} requires each can placed in the box matching its color.}
\label{tab:simtasks}
\end{table}

\paragraph{LIBERO single-arm extension.}
The LIBERO-10 tasks are single-arm with a $7$-D action space and image-only
observation that comprises robot proprioception ($8$-D) plus two RGB views (an agent view and a
wrist view, \texttt{robot0\_eye\_in\_hand}), and there is no privileged state. Episodes are
capped at $520$ steps. We report \textsc{Task~8} (instruction: \emph{put both moka
pots on the stove}) and \textsc{Task~6} (instruction: \emph{put the white mug on the
plate and put the chocolate pudding to the right of the plate}).

\subsection{B.2\quad Base Policies and Frozen Features}

On the dual-arm tasks the base policy is ACT~\cite{zhao2023learning}, finetuned
\emph{separately per task} on that task's demonstrations and frozen thereafter; on
LIBERO the base is a served $\pi_0$ VLA~\cite{black2024pi0}, likewise frozen. Both
the residual's actor/critic and the value's state representation read the base's
\emph{frozen} encoder feature $\psi$, never a privileged object pose. For ACT,
$\psi$ is the mean-pooled encoder tokens (width $512$) concatenated with
proprioception, giving a per-task input of $530$ dims (\textsc{ThreePieceAssembly},
\textsc{Threading}; $512{+}18$), $548$ (\textsc{Pouring}; $512{+}36$), or $550$
(\textsc{LiftTray}; $512{+}38$). For $\pi_0$, $\psi$ is the mean-pooled prefix
feature concatenated with standardized proprioception. \todo{report the $\pi_0$
prefix-feature width (produced server-side at cache-build time), and confirm the
LIBERO proprio dimension used.} The base action chunk is executed open-loop with an
execution horizon of $10$ steps (a $30$-step variant is also run on LIBERO, drawn
solid/dashed in the main paper's LIBERO figure). Ablation~D additionally trains the
value on a privileged $30$-D end-effector\,$+$\,object state (\texttt{eef\_piece}) as
an upper-bound reference; the deployed representation is always image-derived.

\subsection{B.3\quad Configuration and Hyperparameters}

The goal-conditioned value and navigator $\pi_h$ are initialized once per task from
demonstrations and frozen features $\psi$, then jointly finetuned during residual RL in
the waypoint-enabled arms while the bottleneck $\phi$ remains frozen. The separate
single-state value supplying $\Phi_V$ is pretrained with the goal marginalized out and
remains frozen throughout residual training. Table~\ref{tab:simvalue} lists
the value/navigator configuration and Table~\ref{tab:simresrl} the residual-RL
configuration, whose venue-independent core is the shared set of Table~\ref{tab:sharedhp}. Every setting below is shared across the four long-horizon simulated
tasks; the waypoint horizon differs between the dual-arm suite ($15$) and LIBERO ($25$).
Section~C.5 separately states which optimizer settings are reused and which collection
parameters change in imagination.
The concrete values below instantiate the general update of Section~A; the
corresponding training loop is Algorithm~\ref{alg:shore}, instantiated with the simulator
collector and the simulation replay sampler.

\begin{table}[t]
\centering\small\setlength{\tabcolsep}{4pt}
\begin{tabular}{@{}ll@{}}
\toprule
Setting & Value \\
\midrule
Goal-cond.\ value $V(s,\phi)$        & $3\times512$, LayerNorm$+$GELU \\
\quad critic form                    & double, target $\min(V_1,V_2)$ \\
Waypoint bottleneck dim $z$           & $10$ (normalized to radius $\sqrt{10}$) \\
Navigator $\pi_h$                    & $2$ layers $\times$ $256$ \\
\quad AWR inverse temp.\ $\beta$      & $1.0$ (weight clipped at $100$) \\
Waypoint horizon $k$ (way steps)      & $15$ dual-arm / $25$ LIBERO \\
Expectile $\tau$                     & $0.7$ \\
Discount $\gamma$                    & $0.99$ \\
Target EMA                           & $0.005$ \\
Goal sampling                        & geometric($\gamma$), rand.-goal $0.3$ \\
Goal-reach reward (gc)               & $0$ at goal, $-1$ otherwise \\
Single-state $V(s)$ reward           & $1$ at final step, $0$ otherwise \\
Optimizer / lr                       & Adam / $3{\times}10^{-4}$ \\
Steps / batch                        & $50{,}000$ / $256$ \\
\bottomrule
\end{tabular}
\caption{Offline goal-conditioned value and navigator configuration (also the
recipe for the single-state potential value, with a state-only input). Widths
correspond to the \texttt{hiqlv512} run recipe; \texttt{sg15} sets the waypoint
horizon to $15$.}
\label{tab:simvalue}
\end{table}

\begin{table}[t]
\centering\small\setlength{\tabcolsep}{4pt}
\begin{tabular}{@{}p{0.40\columnwidth}p{0.54\columnwidth}@{}}
\toprule
Setting & Value \\
\midrule
Regime                          & RLPD, TD3-style updates \\
Critic ensemble                 & $10$ Q-nets; target $=\min$ of $2$ random heads \\
Critic / actor MLP              & $2$ layers $\times$ $1024$, ReLU, LayerNorm \\
Residual actor output           & $\tanh \times s_a$; final layer zero-init \\
Action scale $s_a$              & $0.05$ \\
Discount $\gamma$ / $n$-step    & $0.99$ / $3$ \\
Update-to-data (UTD)            & $4$ (actor updated every $4$th critic step) \\
Target Polyak $\tau$            & $0.005$ \\
Exploration noise               & $0.05$; target smoothing $0.05$, clip $0.3$ \\
actor\_lr / critic\_lr          & $1{\times}10^{-6}$ / $1{\times}10^{-4}$ (AdamW) \\
Batch / replay buffer           & $256$ / $2{\times}10^{5}$ \\
Offline demo fraction           & $0.5$ (RLPD demo mixing) \\
Demo-BC coefficient             & $0.1$ (MSE in residual coordinates; TODO verify scale) \\
Learning starts                 & $10^{4}$ \\
Total env.\ steps               & $5{\times}10^{5}$ \\
Potential $\Phi=c\,V(s)$        & $c=\frac{\#\text{stages}-1}{V_{\max}-V_{\min}}$, mult.\ $1.0$ \\
Stage bonus (main)              & amplitude $1.0$, stage-balanced replay \\
Gradient clip                   & $1.0$ \\
Eval                            & $50$ ep.\ (DexMimicGen) / $10$ (LIBERO) \\
\quad                          & $8$ envs, deterministic ($\sigma{=}0$) \\
\bottomrule
\end{tabular}
\caption{Residual off-policy RL configuration (simulation). Despite the ``TD3''
label in the main text, the critic is an RLPD-style $10$-head ensemble with
TD3-style target smoothing and delayed actor updates; the residual actor's final
layer is zero-initialized, so the residual is exactly $0$ at step $0$ and the
leftmost evaluation of every run measures the pure frozen base (the BC floor;
Section~B.4).}
\label{tab:simresrl}
\end{table}

\subsection{B.4\quad Baselines}

We compare against four external references and our own component ablations. A comparison
is described as \emph{matched} only when it uses the common base, data, budget, execution
horizon, evaluation harness, and seed policy specified by the evaluation protocol (Section~B.6).
\emph{BC base} is the frozen base policy with no RL: because the residual's actor is
zero-initialized (Section~B.3), it is read for free as the leftmost evaluation point
($\text{env steps}\approx 0$) of every run, evaluated over the same $50$ episodes and
the same success predicate as every later point. \emph{Flat demo-anchored
residual}~\cite{resfit2025residual} is the published residual recipe (our arm~$2$,
below); the \emph{naive} published recipe without the demo anchor (arm~$1$) is the
one that collapses. \emph{DSRL}~\cite{wagenmaker2025dsrl} places the residual in the
base's noise space rather than raw action space, on the \emph{same} base.
\emph{IBRL}~\cite{hu2023imitation} is a BC-bootstrapped off-policy method that
soft-anchors via a $Q$-argmax over base and RL actions. \emph{Full-policy
IQL}~\cite{kostrikov2022iql} finetunes the whole policy with no frozen base,
isolating the residual structure itself.

\subsection{B.5\quad Ablation Design}

All ablations read off a single shared set of seven arms
(Table~\ref{tab:simarms}), run inside the \emph{same} codebase on identical frozen
features, offline data, action scale, seed set, and budget, so each comparison
isolates one factor. Online joint finetuning of the value head and navigator is
bundled with waypoint conditioning (enabled whenever a waypoint is injected), so
``waypoint'' denotes waypoint$+$joint as one mechanism. The potential takes one of
three forms, namely none, the hand-designed per-stage bonus (\emph{stage}, the form used by
the main result), or the self-derived single-state $V(s)$.

\begin{table}[t]
\centering\small\setlength{\tabcolsep}{4pt}
\begin{tabular}{@{}llccc@{}}
\toprule
\# & Arm & waypoint($+$joint) & potential & BC \\
\midrule
1 & flat, BC$0$             & off & none      & $0$   \\
2 & flat, BC$0.1$           & off & none      & $0.1$ \\
3 & waypoint, BC$0$          & on  & none      & $0$   \\
4 & waypoint, BC$0.1$        & on  & none      & $0.1$ \\
5 & potential-only          & off & $V(s)$    & $0.1$ \\
6 & self-derived pot.       & on  & $V(s)$    & $0.1$ \\
7 & \textbf{SHORE-RL (+stage)} & on & stage   & $0.1$ \\
\bottomrule
\end{tabular}
\caption{Shared ablation arms (per task; fixed action scale/seed/budget). Arm~$7$ is
the headline configuration; arm~$6$ is the same thing with the hand-designed bonus
replaced by the self-derived $V(s)$. Arms $1$--$4$ form the BC$\times$waypoint
$2\times2$ (Ablation~A); arms $2$,$4$,$5$,$6$ the waypoint$\times$potential $2\times2$
(Ablation~B); arms $7$,$6$ the hand-designed-vs-self-derived pair (Ablation~C, on
stage-structured tasks only). Ablation~D varies the representation $\psi$
(privileged \texttt{eef\_piece} vs.\ ACT vs.\ $\pi_0$ features) and is orthogonal to
this table.}
\label{tab:simarms}
\end{table}

\paragraph{Coverage.}
So that no conclusion rests on a single task, the full seven-arm matrix is run on
\textsc{Pouring} ($3$ seeds), and the decisive arms (flat$+$BC, waypoint-only,
potential-only, SHORE-RL, and the self-derived $V(s)$ substitution) are replicated on
the other three long-horizon tasks (\textsc{ThreePieceAssembly}, \textsc{Threading},
\textsc{LiftTray}; $1$--$2$ seeds). The short-horizon control \textsc{CanSort} is
excluded because no arm collapses there.

\subsection{B.6\quad Evaluation Protocol}

All arms are run under one evaluation harness so that no method is advantaged by its own
setup. The intended matched protocol gives all arms the same base, the same offline
demonstrations, the same evaluation harness ($50$-episode deterministic evaluation for
DexMimicGen, $10$ for LIBERO; the same execution horizon and success predicate), the same
budget ($5{\times}10^{5}$ env steps), and the same seed set (${\geq}3$ on the flagship
task). The respective anchoring knobs, namely our action scale, DSRL's noise scale, and
IBRL's exploration $\sigma$, must be calibrated under that harness before a strict
comparison is reported. \todo{run DSRL and IBRL under this matched harness on the dual-arm
tasks; the LIBERO DSRL overlay in the main paper uses DSRL's own harness ($10$ ep./eval,
base re-queried every $20$ steps) and is therefore indicative rather than strictly
controlled.}

\paragraph{Reported metric.}
The reported metric is the \emph{steady-state} success rate, defined as the mean over the
last $20\%$ of evaluations, each of $50$ deterministic episodes (\textsc{CanSort} and the
DexMimicGen tasks) or $10$ (LIBERO). \todo{define whether the last-$20\%$ mean is computed
within each seed before aggregating, report dispersion across seeds, and predefine how
incomplete or crashed runs are excluded.}

\subsection{B.7\quad Results, Run Inventory, and Provenance}

The main paper reports the compact headline table; this subsection should provide the
reader-facing per-task, per-arm, and per-seed values underlying it, together with the run
inventory needed to interpret incomplete cells. The reported metric is the steady-state
success rate defined in Section~B.6.
\todo{add the full per-seed result table and learning curves.}

\paragraph{Run inventory.}
\textsc{Pouring} carries the full seven-arm matrix at $3$ seeds to
$5{\times}10^{5}$ steps. On \textsc{ThreePieceAssembly}, \textsc{Threading}, and
\textsc{LiftTray} the decisive arms are run at $1$--$2$ seeds; the hand-designed
\emph{stage} arm reaches $5{\times}10^{5}$ at $3$ seeds where reported. The
self-derived $V(s)$ arm reached $5{\times}10^{5}$ (single seed) on
\textsc{LiftTray} and \textsc{Pouring} and holds there, but crashed early on a shared
machine on \textsc{ThreePieceAssembly} (${\sim}2{\times}10^{4}$) and
\textsc{Threading} ($2.4$--$2.6{\times}10^{5}$), which was an infrastructure failure rather than a
performance collapse, so its late-window behavior is not yet observed on those two
tasks. \todo{re-run the crashed $V(s)$ arms to $5{\times}10^{5}$ and add seeds;
report per task, not pooled.}

\paragraph{Provenance items to resolve before the result table is final.}
(i)~The \textsc{Pouring} \emph{flat residual} cell must be traced to a specific run:
its value contradicts the collapsed \textsc{Pouring} baseline curve, and whether that
row carries the demo anchor (BC$0.1$) or not changes whether the anti-collapse claim
is attributed to the BC anchor or the waypoint.
(ii)~The \textsc{Threading} \emph{$+$stage} cell had a stale entry; read the final
number off the $3$-seed $5{\times}10^{5}$ run.
(iii)~The BC-base cells use the step-$0$ protocol above; the frozen-base reference
shifted when the base runs switched to execution horizon $10$, so any pre-switch base
number must not be reused. \todo{resolve (i)--(iii) and fill the DSRL/IBRL/IQL rows
from the matched-harness runs of Section~B.6.}
\todo{after resolving the items above, replace this internal provenance checklist with a
compact reader-facing table containing the run identifier, seed count, final step, and
inclusion status for every reported cell.}

\subsection{B.8\quad Limitations of the Simulation Venue}

\paragraph{Seeds and unfinished runs.}
Several cells carry $1$--$2$ seeds against the flagship $3$, and the self-derived
$V(s)$ arm has not been observed through the late window on two of four tasks
(Section~B.7). The parity claim between the self-derived $V(s)$ and the hand-designed
stage bonus is therefore suggestive, not yet conclusive, and is stated as such.

\paragraph{The stage bonus is privileged.}
The main result's per-stage bonus reads a task-specific simulator predicate and only
exists on tasks with a clean stage decomposition; it is exactly this dependence that
the self-derived $V(s)$ is meant to remove, and it is what makes the real-robot venue
(Section~C) use the $V(s)$ form throughout.

\paragraph{Cross-method comparison on LIBERO.}
The LIBERO DSRL overlay runs under DSRL's own evaluation harness rather than ours, so
absolute-success comparisons there are indicative rather than strictly controlled
(Section~B.6).

\section{C\quad Real-Robot Venue}

Everything in Section~B trains the residual against a cheap-to-step simulator. For
hardware evaluation we retain the residual parameterization and actor--critic objective of
Section~A, but instantiate its collector and replay sampler with a learned dynamics model
finetuned on our own robot data. This section documents that venue: the platform and its
interfaces (C.1), the three tasks and their success
criteria (C.2), the demonstration and rollout data (C.3), the dynamics model (C.4),
the imagination loop the residual is trained in (C.5), the protocol used to evaluate
the resulting residual on the physical robot (C.6), and the limitations specific to
this venue (C.7).

\subsection{C.1\quad Real-Robot Platform and Control}

\paragraph{Robot.}
The real-robot venue uses a TeleAvatar bimanual platform: two $7$-DoF arms, each
terminated by a parallel gripper. Commands are joint-space targets, so the action
space is $16$-dimensional, $(7\ \text{joints} + 1\ \text{gripper}) \times 2$.
\todo{confirm the low-level controller rate and whether commands are absolute
joint positions or deltas.}

\paragraph{Perception.}
Perception is image-only: three RGB cameras stream $848{\times}480$ video at
$30$\,Hz, consisting of one head-mounted view (\texttt{top\_head}) and two wrist-mounted views
(\texttt{hand\_left}, \texttt{hand\_right}). The same three views are consumed by
every component that needs to see the scene: the base policies, the frozen
encoder $\psi$ that supplies the value's state representation, and the dynamics
model. No component receives a signal the deployed robot does not have, because there is
no motion capture, no object pose, and no privileged state anywhere in this
venue.

\paragraph{Base policy and action chunking.}
Each task's base policy is a served $\pi_{0.5}$-style VLA~\cite{sima2026kai0} that
emits $50$-step
action chunks, executed open-loop within each chunk. Internally the VLA's action
vector is zero-padded to its native width of $32$ dimensions, of which the
leading $16$ are the platform's commands; the padding is stripped before
execution and before any residual is added.

\paragraph{What runs at deployment.}
At deployment the robot runs the task's frozen base policy plus the small
residual and the waypoint head, both evaluated on base features. The dynamics
model is a training-time device only and is never executed at inference time, so
the deployed system carries no video model and no additional latency beyond the
two small heads.

\subsection{C.2\quad Tasks and Success Criteria}

We evaluate on three tabletop bimanual tasks. Episodes run $800$--$2{,}500$ steps
at $30$\,Hz ($\approx 27$--$83$\,s). For each task we quote the language
instruction verbatim as it appears in the dataset, because that same string is
the caption conditioning the dynamics model (Section~C.4), and because a mismatch
between the two would silently change what the model is conditioned on.
\todo{the \textsc{pick\_paper\_roll} data is recorded at $20$\,Hz (Section~C.3,
Table~\ref{tab:realdata}), so its episode-length and duration range differ from
the $30$\,Hz figure quoted above; reconcile this range with the per-task frame
counts before submission.}

\paragraph{\textsc{pick\_paper\_roll}} (instruction: \texttt{put the paper roll on
the holder}). The robot grasps a paper roll from the tabletop and places it onto
an upright holder. A key challenge is that the roll is a rolling object: a grasp
that is not centered displaces it rather than securing it, and the resulting
error compounds into the placement. \emph{Success}: the paper roll is seated on
the holder and the gripper has opened and withdrawn, with the roll remaining on
the holder. \todo{confirm whether a handover between arms is required, and the
exact seated/inserted criterion.}

\paragraph{\textsc{pick\_cup}} (instruction: \texttt{pick cup}). The robot grasps
a cup and places it at the target location. \emph{Success}: the cup is upright at
the target and the gripper has released it, with the cup remaining upright.
\todo{confirm the target (plate/coaster/tray?), the cup's initial orientation
(upright or lying down), and whether reorientation is part of the task.}

\paragraph{\textsc{build\_block}} (instruction: \texttt{build block}). An
assembly-style task in which blocks are stacked. \emph{Success}: the blocks are
stacked in the required configuration and the stack remains standing after the
gripper releases. \todo{confirm the number of blocks, the required stacking
order, and whether each arm is assigned particular blocks.}

\paragraph{Randomization.}
Object placements vary across episodes. \todo{confirm the randomization protocol:
which objects are randomized, over what region, whether the robot's initial pose
is randomized, and whether the evaluation initial-state set (Section~C.6) is drawn
from the same distribution as the training placements or held out.}

\paragraph{Why binary success.}
We score each trial with a single binary criterion rather than a per-stage
rubric. This keeps the real-robot metric identical in kind to the simulated one
and avoids introducing stage labels into a venue whose entire point is that it
runs without them: the shaping potential here is the self-derived single-state
value, precisely because no stage detector is available on hardware.

\subsection{C.3\quad Demonstration and Rollout Data}

Two kinds of data are collected on the physical robot, and they enter the
pipeline at different points.

\paragraph{Expert demonstrations.}
Human teleoperation, success-only, stored in LeRobot format with per-frame
images, joint states, and actions. These train the base policies: we finetune the
$\pi_{0.5}$ VLA \emph{separately per task}, giving one base policy per task, each
frozen from that point on.

\paragraph{Base-policy rollouts.}
We then deploy each frozen base on the physical robot and record its rollouts.
These are mixed-outcome by construction, in that they contain exactly the residual
failures that the base policy has not overcome, and they are used only for the
dynamics model, never for base-policy training. This is deliberate: the model
must be able to render the failure modes the residual is being trained to avoid,
and a model finetuned on success-only data cannot, because it has never seen one.

\begin{table}[t]
\centering\small\setlength{\tabcolsep}{4pt}
\begin{tabular}{@{}lrrrrr@{}}
\toprule
Task & expert & rollout & total & frames & Hz \\
\midrule
\textsc{build\_block}      & $295$ & $50$ & $345$ & $378{,}724$ & $30$ \\
\textsc{pick\_cup}         & $207$ & $51$ & $258$ & $233{,}589$ & $30$ \\
\textsc{pick\_paper\_roll} & $849$ & $53$ & $902$ & $289{,}278$ & $20$ \\
\midrule
total                      & $1{,}351$ & $154$ & $1{,}505$ & $901{,}591$ & mixed \\
\bottomrule
\end{tabular}
\caption{Real-robot data, per task. Each task is one dynamics-model domain
(\texttt{block}, \texttt{cup}, \texttt{paper} respectively; Section~C.4.2).
\emph{Expert} are success-only human
teleoperation episodes and train that task's base policy; \emph{rollout} are
recorded deployments of the resulting frozen base policy and enter only the
dynamics-model finetuning (Section~C.4). Frame counts are for the merged
per-task dataset. The \textsc{pick\_paper\_roll} data is recorded at $20$\,Hz
rather than $30$\,Hz; Section~C.4.2 explains why this is handled within the domain
rather than across domains. \todo{the main paper states ${\sim}300$ success-only
episodes per task, which matches \textsc{build\_block} and \textsc{pick\_cup} but
not \textsc{pick\_paper\_roll} ($849$); reconcile the two before submission.}}
\label{tab:realdata}
\end{table}

\subsection{C.4\quad Dynamics Model}

\subsubsection{C.4.1\quad Architecture and Action Conditioning}

The dynamics model $D$ predicts the next $H$ multi-view frames from an
observation history and an action chunk. We instantiate it as an
action-conditioned multi-view video-diffusion model and finetune it from a
publicly released checkpoint pretrained on large-scale robot-manipulation
video~\cite{yang2026rise}, following that work's training recipe.

The backbone is an LTX-Video stack, comprising a T5 text encoder, a causal video VAE, and a
$28$-layer diffusion transformer ($32$ attention heads of dimension $64$,
cross-attention width $2{,}048$, $128$ latent channels), and it is trained with a flow-matching
objective under logit-normal timestep weighting. Each training sample is
$3$ views $\times$ $[\,4\ \text{history} + 25\ \text{predicted}\,]$ frames at
$256{\times}192$, downsampled from the native $848{\times}480$, paired with the
$50$-step action chunk that spans the predicted window.

Actions condition the model through a dedicated path: the chunk is normalized
with per-domain statistics, zero-padded from $16$ to the backbone's $30$-channel
action width, encoded by a small ViT, and \emph{added} to the caption embeddings
that drive cross-attention. This is the pretrained checkpoint's own conditioning
path, so no part of it is randomly initialized. It is, however, a comparatively
weak form of conditioning, being an additive term on a text embedding rather than a
separate control stream, which is why action controllability is something we
test rather than assume (Section~C.4.4), and why we flag it in Section~C.7.

\subsubsection{C.4.2\quad Mixed-Quality Training Data and Shortcut Controls}

A \emph{single} dynamics model is jointly full-finetuned on the pooled data of
all three tasks. The asymmetry with the base policies, which are one per task against a single
shared dynamics model, is deliberate: predicting how a scene responds to an
action chunk is largely task-independent and benefits from pooling, whereas the
policies do not.

The data is organized as three domains, one per task, each containing that task's
expert demonstrations \emph{and} that task's base-policy rollouts merged into a
single dataset (Table~\ref{tab:realdata}). Because the model has no notion of
success or failure, its objective being pure video prediction, mixing the two is
what teaches it to render failure at all.

That mixing, however, creates an opportunity for the model to learn a shortcut
instead of dynamics, and we control for two such shortcuts explicitly.

\paragraph{(i) Frame rate, within a domain.}
The \textsc{pick\_paper\_roll} expert data is recorded at $20$\,Hz while its
failure rollouts were recorded at $30$\,Hz. The loader samples windows by
\emph{frame count}, not by wall-clock duration, and the model is never told the
frame rate, since its temporal position encoding is the bare frame index. A
within-domain frame-rate difference would therefore make inter-frame motion
magnitude collinear with success versus failure, letting the model infer the
outcome from apparent speed rather than from the action-conditioned dynamics,
which is exactly the capability the negative probe of Section~C.4.4 is meant to
test. We resample the failure rollouts to $20$\,Hz so that the domain is
internally uniform. The remaining \emph{across}-domain difference
($20$\,Hz for \texttt{paper} versus $30$\,Hz for \texttt{block} and
\texttt{cup}) is left as is: it is collinear with the task, not with the
outcome, and task identity is already supplied to the model through the caption.

\paragraph{(ii) Caption, within a domain.}
For the same reason, all episodes within a domain carry one identical caption, so
text cannot serve as a success/failure indicator. This required unifying the
\texttt{block} domain, whose expert and rollout splits had been captioned
differently, onto the single string \texttt{build block}.

\subsubsection{C.4.3\quad Finetuning Configuration}

Finetuning is a single joint, full-parameter update over the pooled data of all
three domains. Table~\ref{tab:wmhp} lists the configuration.

\begin{table}[t]
\centering\small\setlength{\tabcolsep}{4pt}
\begin{tabular}{@{}ll@{}}
\toprule
Setting & Value \\
\midrule
Initialization      & released pretrained checkpoint \\
Objective           & flow matching (logit-normal weighting) \\
Views               & $3$ (head, left wrist, right wrist) \\
Frames per sample   & $4$ history $+$ $25$ predicted \\
Resolution          & $256{\times}192$ \\
Action chunk        & $50$ steps, per-domain normalized \\
Optimizer           & AdamW, $\beta=(0.9, 0.95)$ \\
Learning rate       & $1{\times}10^{-4}$, constant \\
Warmup steps        & $2{,}000$ \\
LR scaling by batch & none \\
Precision           & bf16, gradient checkpointing \\
Sharding            & DeepSpeed ZeRO-2 \\
Batch size          & $4$ per GPU ($48$\,GB) / $8$ per GPU ($96$\,GB) \\
Training steps      & $20{,}000$ \\
Caption dropout     & $0.06$ \\
First-frame noise   & $0.2$ \\
Validation          & every $1{,}000$ steps (per-domain videos) \\
Checkpoint          & every $2{,}000$ steps \\
\bottomrule
\end{tabular}
\caption{Dynamics-model finetuning configuration. \todo{P0: this is the verified
configuration, and the pipeline has been validated end-to-end at $30$ steps
(loss $0.107$; the pretrained checkpoint loads with no missing or unexpected
tensors), but the full $20{,}000$-step run must be confirmed complete, and its
wall-clock and hardware reported, before this table is described as the
configuration of a trained model.}}
\label{tab:wmhp}
\end{table}

\subsubsection{C.4.4\quad Validating the Model Before Any Policy Training}

The dynamics model is the environment in this venue, so its failure modes become
the residual's failure modes. We therefore validate it on held-out episodes
before training any policy against it, along three axes; these are the three
panels of the validation figure in the main paper.

\paragraph{(i) Imagined-rollout quality.}
Given a real observation history and the ground-truth action chunk that follows
it, the imagined $H$-step future is compared against the recorded future, and we check whether
object position and grasp timing line up.

\paragraph{(ii) Action controllability, including a negative probe.}
The same observation history under different action chunks must yield
\emph{diverging} futures; a model that ignores its action input is a video prior,
not a dynamics model, and would make the residual's gradient meaningless. The
sharper form of this test is the negative probe: we feed deliberately perturbed
actions, such as translating the grasp segment by a few centimeters or closing the gripper
early, and require the imagined video to show failure (a missed grasp, a slip)
rather than a hallucinated success. Passing this probe is the entire reason the
mixed-quality data of Section~C.4.2 exists; failing it means the failure fraction
or the training budget is insufficient.

\paragraph{(iii) Potential-trajectory consistency.}
Finally we check the quantity the residual actually optimizes, rather than
pixels: $V(\psi(\hat{o}_{t+k}))$ must track $V(\psi(o_{t+k}))$ over $k$. A model
can be visually plausible and still be useless here if the re-encoded imagined
frames land where the potential misreads them, so this is the axis that most
directly predicts whether imagination training transfers.

\todo{P0: run all three and report; state the held-out split size and, for (i),
the metric used.}

\subsection{C.5\quad Residual RL in Imagination}

\paragraph{Why a replay buffer rather than a step function.}
A video-diffusion model takes seconds per generation and cannot serve as a
per-step simulator inside an off-policy loop. We therefore decouple generation
from updating and train from an \emph{imagination replay buffer}: every $K$
policy updates we draw $M$ starting states from the real-robot data, roll $D$
forward $H$ steps under the residual-perturbed actions
$a = a_b + s_a\,a_{\text{res}}$, score the imagined frames once with the
potential $\Phi$ and cache the result, and then let the
TD3~\cite{fujimoto2018addressing} residual perform many
cheap updates on a mixture of real and imagined transitions. Imagination is
recursed at most twice from any real state, bounding how far model error can
compound before the trajectory is re-anchored to real data.
\todo{fill $K$, $M$, $H$, the demonstration/recorded/imagined mixing ratio, the
update-to-data ratio, buffer capacities, learning-start threshold, and wall-clock per
$10^5$ updates; define which recorded transitions receive which reward and how their
terminal flags are constructed.}
This loop is the imagination instantiation of Algorithm~\ref{alg:shore}, with
$\mathcal C$ set to bounded rollouts of $D$ and $\rho$ set to the reported replay mixture.

\paragraph{The reward is the potential difference.}
Imagined rollouts carry no environment success signal, and the hand-designed
per-stage bonus used in the main simulated results (Section~B.3) reads simulator
predicates that do not exist here. The per-step reward in this venue is therefore the
potential difference alone, $\tilde r_t=\gamma\Phi_V(s_{t+1})-\Phi_V(s_t)$, with
$\Phi_V=cV(s)$ evaluated on re-encoded imagined frames. No frame scorer and no reward
model is introduced.
The safety rule of the main paper applies unchanged: shaping uses only the
single-state value, while the goal-conditioned value contributes only the waypoint
$z$ and never the potential.

\paragraph{What is frozen and what is trained.}
Frozen: the base policy, the feature encoder $\psi$, the dynamics model $D$, and
the single-state value supplying $\Phi$. Trained: the residual actor and critic
and, as in simulation, the goal-conditioned value head and the navigator
$\pi_h$, jointly finetuned online while $\phi$ stays frozen. The actor--critic settings are the venue-independent ones of
Table~\ref{tab:sharedhp} (Section~A); collection and replay settings are venue-specific
and must be reported above.

\subsection{C.6\quad Real-Robot Evaluation Protocol}

The residual trained in imagination is evaluated on the physical robot. Since
the venue changes but the algorithm does not, the comparison is the same one the
simulated experiments make: the task's frozen base policy versus base $+$
SHORE-RL residual. In each task the two conditions deploy the identical frozen base
and identical controller and differ only in whether the residual is active.

\paragraph{Trials.}
$50$ trials per cell, drawn from a fixed initial-state set.

\paragraph{Paired initial states.}
The two conditions are evaluated from the \emph{same} physical initial configurations,
so the comparison is paired rather than between independent samples. Tabletop
object poses are restored from a marked template between trials; the pairing is
what makes a $50$-trial difference interpretable at all, since initial-state
difficulty on these tasks varies more than the effect we are trying to measure.

\paragraph{Interleaving.}
The two conditions alternate within a session rather than being run in blocks, so that
session-level drift, such as lighting, gripper wear, or calibration, cannot be confounded
with the condition.

\paragraph{Blind scoring.}
Trials are recorded and scored offline against the Section~C.2 criteria by an
annotator who is not told which condition produced the trial.

\paragraph{Reporting.}
We report the mean with its standard error, together with Wilson score $95\%$
intervals for a binomial with $n=50$, which is better behaved than the normal
approximation near the extremes where several of these cells sit. Because trials are
paired, the primary comparison should additionally report the paired success difference,
the two discordant-pair counts, and a paired confidence interval or exact McNemar test.
\todo{pre-specify the paired statistical analysis and the treatment of aborted or invalid
trials before unblinding condition labels.} We additionally report, per task, a breakdown
of how many of the $50$ trials failed and in what
manner (Table~\ref{tab:failmodes}), since an aggregate success rate cannot
distinguish a residual that fixes the base's dominant failure from one that
trades it for another.

\subsection{C.7\quad Real-Robot Results and Failure Analysis}

This subsection reports physical-robot performance only after the dynamics model passes
all three validation axes of Section~C.4.4 and the paired protocol above is complete.

\begin{table}[t]
\centering\small\setlength{\tabcolsep}{4pt}
\begin{tabular}{@{}lccc@{}}
\toprule
Task & base & \textbf{SHORE-RL} & paired $\Delta$ \\
\midrule
\textsc{pick\_paper\_roll} & \todo{P1} & \todo{P1} & \todo{P1} \\
\textsc{pick\_cup}         & \todo{P1} & \todo{P1} & \todo{P1} \\
\textsc{build\_block}      & \todo{P1} & \todo{P1} & \todo{P1} \\
\bottomrule
\end{tabular}
\caption{Physical-robot success over paired initial states. Report successes out of $50$,
success rate with Wilson $95\%$ interval for each arm, paired difference with its interval,
and the two discordant-pair counts in the text.}
\label{tab:realresults}
\end{table}

\todo{P1: fill Table~\ref{tab:realresults}; report the paired test, session count and dates,
invalid-trial count, and whether all tasks used the same pre-registered initial-state set.}

\begin{table}[t]
\centering\small\setlength{\tabcolsep}{4pt}
\begin{tabular}{@{}llcc@{}}
\toprule
Task & Failure mode & base & \textbf{SHORE-RL} \\
\midrule
\textsc{pick\_paper\_roll} & \multicolumn{3}{c}{\todo{P1}} \\
\textsc{pick\_cup}         & \multicolumn{3}{c}{\todo{P1}} \\
\textsc{build\_block}      & \multicolumn{3}{c}{\todo{P1}} \\
\bottomrule
\end{tabular}
\caption{Failure-mode breakdown, counts out of $50$ trials per cell.
\todo{P1: define the failure-mode categories per task from the recorded trials
(e.g.\ missed grasp, slip during transport, misplacement), then fill. Categories
must be fixed before scoring, not after.}}
\label{tab:failmodes}
\end{table}

\subsection{C.8\quad Limitations of the Real-Robot Venue}

\paragraph{The environment is a model.}
The residual is optimized against $D$, not against the robot, so any systematic
error in $D$ becomes a systematic bias in the learned residual. We bound the
compounding with a recursion depth of two and test $D$ before use
(Section~C.4.4), but neither eliminates the gap; the honest reading of this venue
is that it trades an unbiased but unaffordable environment for a biased but
cheap one.

\paragraph{No environment reward to correct the potential.}
In simulation the shaped reward sits alongside a real terminal success signal. In
imagination it does not: the reward is derived entirely from changes in the learned
potential. A
mis-calibrated $V(s)$ therefore has nothing to correct it, which makes axis
(iii) of Section~C.4.4 load-bearing rather than a nicety.

\paragraph{Weak action conditioning.}
As noted in Section~C.4.1, actions enter $D$ additively on the caption embedding.
Whether perturbed actions reliably render failure is an empirical question the
negative probe has to answer for each task; a task where it does not is a task
where this venue does not apply.

\paragraph{One model, three tasks.}
Pooling buys data efficiency but means the three tasks share capacity; we do not
currently measure whether per-task models would predict better.

\paragraph{Real-robot interaction is evaluation-only.}
No RL touches the hardware: the robot is used to collect demonstrations, to
record base-policy rollouts, and to evaluate. This is what makes the venue safe
and reset-free, but it also means we never observe how the residual would behave
under genuine on-robot exploration.

% Supplementary Content Ends Here

% References and End of Paper
% These lines must be placed at the end of your paper
\bibliography{references}

\end{document}
